\section{Experimental Setup}

\textbf{Dataset and partitioning.}
All reported experiments use CIFAR-10. Near-IID settings use a large Dirichlet
concentration ($\alpha=100$), while the heterogeneous setting uses
$\alpha=0.5$.

\textbf{Model and local training.}
Clients train a shared CNN for two local epochs per communication round. The
server-side CorrectionNet uses hidden width 256.

\textbf{Client types.}
Each experiment uses 10 total clients with full participation. Anchor fractions
range from 20\% to 50\%, depending on the configuration.

\textbf{Baselines.}
We compare against FedAvg \citep{mcmahan2017} and a simplified ProWD-style
baseline derived from \citet{yoon2022prowd}. The simplified ProWD baseline
implemented in this repository uses a shared three-layer residual MLP trained
only on the current round and targets the mean anchor delta.

\textbf{Reported configurations.}
Table~\ref{tab:configs} summarizes the main settings referenced in the paper and
slides.

\begin{table}[h]
  \centering
  \caption{Reported experimental configurations.}
  \label{tab:configs}
  \small
  \resizebox{0.95\columnwidth}{!}{%
  \begin{tabular}{lccccc}
    \toprule
    Setting & $\alpha$ & Bits & Anchors & Rounds & Corr. epochs \\
    \midrule
    Near-IID ProWD benchmark & 100 & INT8 & 5 / 10 & 50 & 5 \\
    Near-IID severe quant.   & 100 & INT2 & 2 / 10 & 50 & 20 \\
    Near-IID moderate quant. & 100 & INT4 & 5 / 10 & 50 & 10 \\
    Non-IID heterogeneous    & 0.5 & INT8 & 5 / 10 & 50 & 5 \\
    \bottomrule
  \end{tabular}}
\end{table}

This table is mainly an interpretation guide for the later accuracy results.
The near-IID INT8 row is the easiest setting and is closest to the regime where
plain averaging is already strong. The INT2 row is much harsher: only two
anchors supervise eight constrained clients, and each coordinate is compressed
to just three effective levels. Because the rows differ in both bit-width and
anchor fraction, the most important comparisons are within each row across
methods, while the cross-row trend tells us when correction becomes necessary.

\textbf{Evaluation.}
The primary metric is top-1 accuracy on the global CIFAR-10 test set at
selected communication rounds. We focus on round 10 to measure early usefulness
and round 50 to measure whether gains persist.
